\par\Needspace{6\baselineskip}\medskip
\begin{center}\hypertarget{native-183-article137}{}\pdfbookmark[2]{137.}{article-137}137.\end{center}
The \emph{second} task is to determine the relative orientation of those three great circles.
Their relative orientation depends on the positions of their mutual intersections and on their
inclinations. If we wish to express these relationships unambiguously in clear and general
terms, so that we need no separate figures for the various cases, a few preliminary explanations
are needed. \emph{First}, on every great circle we must distinguish its two opposite
\emph{directions} in some way. We do so by regarding one as forward or positive and the other as
backward or negative. Since this choice is in itself entirely arbitrary, we establish a fixed
rule by always taking the directions from $A,A',A''$ toward $B,B',B''$ as positive. Thus, for example, if
the intersection of the first circle with the second is specified by a positive distance from
the point $A$, the distance is understood to be taken from $A$ toward $B$ (as for $D''$ in
our figure). If the distance were negative, it would have to be taken on the other side of
$A$. \emph{Second}, the two hemispheres into which every great circle divides the whole sphere
must likewise be distinguished by suitable names. We call the hemisphere to the right of an
observer following the great circle in the forward direction on the inner surface of the sphere
the \emph{upper} hemisphere, and the other the \emph{lower} hemisphere. The upper region thus
corresponds to the northern hemisphere relative to the ecliptic or equator, and the lower region
to the southern hemisphere.

With these conventions understood, we can readily distinguish \emph{both} intersections of two
great circles from one another. At one intersection the first circle passes from the lower
region of the second into its upper region, or equivalently the second passes from the upper
region of the first into its lower region. At the other intersection the reverse holds. The
choice of intersections in our problem is in itself entirely arbitrary. To follow a fixed rule
here too, we shall always choose those intersections ($D,D',D''$ in Fig. 4) where, respectively, the
third circle $A''B''$ passes into the upper region of the second $A'B'$, the third into the upper
region of the first $AB$, and the second into the upper region of the first. The positions of
these intersections are determined by \hypertarget{native-184}{}their distances from the points $A'$ and $A''$,
$A$ and $A''$, $A$ and $A'$, which we denote simply by $A'D,A''D,AD',A''D',AD'',A'D''$. The mutual inclinations
of the circles are then the angles at these respective intersections $D,D',D''$ between the portions
of the intersecting circles that lie in the forward direction. These inclinations, always taken
between $0$ and $180^0$, are denoted by $\epsilon,\epsilon',\epsilon''$. Determining these nine unknown quantities from
the known ones plainly depends on the same problem we treated in Article 55. We therefore have
the following equations
\[
\begin{aligned}
3)\quad&\sin\tfrac12\epsilon\sin\tfrac12(A'D+A''D)=\sin\tfrac12(l''-l')\sin\tfrac12(\gamma''+\gamma')\\[.5ex]
4)\quad&\sin\tfrac12\epsilon\cos\tfrac12(A'D+A''D)=\cos\tfrac12(l''-l')\sin\tfrac12(\gamma''-\gamma')\\[.5ex]
5)\quad&\cos\tfrac12\epsilon\sin\tfrac12(A'D-A''D)=\sin\tfrac12(l''-l')\cos\tfrac12(\gamma''+\gamma')\\[.5ex]
6)\quad&\cos\tfrac12\epsilon\cos\tfrac12(A'D-A''D)=\cos\tfrac12(l''-l')\cos\tfrac12(\gamma''-\gamma').
\end{aligned}
\]
Equations 3 and 4 give $\tfrac12(A'D+A''D)$ and $\sin\tfrac12\epsilon$, and the other two give $\tfrac12(A'D-A''D)$ and $\cos\tfrac12\epsilon$. From these
we obtain $A'D,A''D$ and $\epsilon$. The ambiguity in determining the arcs $\tfrac12(A'D+A''D)$, $\tfrac12(A'D-A''D)$ from their
tangents is resolved by the condition that $\sin\tfrac12\epsilon$ and $\cos\tfrac12\epsilon$ must be positive. The agreement
between $\sin\tfrac12\epsilon$ and $\cos\tfrac12\epsilon$ serves to check the whole calculation.

The quantities $AD',A''D',\epsilon',AD'',A'D'',\epsilon''$ are determined in exactly the same way. There is no need to copy out the
eight equations for this calculation, since they follow directly from equations 3---6 by making
the replacements
\begingroup\postdisplaypenalty=10000
\[
\begin{array}{r c|c|c|c|c|c}
 & A'D & A''D & \epsilon & l''-l' & \gamma'' & \gamma'\\
\text{with} & AD' & A''D' & \epsilon' & l''-l & \gamma'' & \gamma\\
\text{or with} & AD'' & A'D'' & \epsilon'' & l'-l & \gamma' & \gamma
\end{array}
\]
respectively.\endgroup

The relations between the sides and angles of the spherical triangle formed by the points $D,D',D''$
provide a further check on the whole calculation. They yield the following equations, which hold
whatever the positions of these points
\[
\frac{\sin(AD'-AD'')}{\sin\epsilon}=\frac{\sin(A'D-A'D'')}{\sin\epsilon'}=\frac{\sin(A''D-A''D')}{\sin\epsilon''}\mathpunct{\cdot}
\]

Finally, if the equator has been chosen \hypertarget{native-185}{}as the fundamental plane in place of the ecliptic,
the calculation changes only in that we must replace the heliocentric positions of the Earth
$A,A',A''$ with the points where the circles $AB,A'B',A''B''$ intersect the equator. Thus for $l,l',l''$ we take
the right ascensions of these intersections, and for $A'D$ the distance of the point $D$
from the second intersection, etc.
